\documentclass[11pt]{article}
\usepackage[margin=1in]{geometry}
\usepackage{amsmath,amssymb,mathtools}
\usepackage{microtype}
\usepackage[colorlinks=true,linkcolor=blue,citecolor=blue,urlcolor=blue]{hyperref}
\newcommand{\Sp}{\mathfrak S}
\newcommand{\R}{\mathbb R}
\title{The Frank--Sabin weak Schatten endpoint:\\
dimension one is false, dimensions $d\ge2$ remain open}
\author{Literature-identification note for arXiv:1609.08388}
\date{12 August 2026}

\begin{document}
\maketitle

\begin{center}
\fbox{\parbox{0.88\linewidth}{\textbf{Status.}
Literature-implied answer (dimension-one counterexample).  The endpoint is
false for $d=1$.  No answer is claimed for $d\ge2$, where a 2025 source still
records the conjecture as open.}}
\end{center}

\section{The source question}

Let
\[
 S=\{(\omega,\xi)\in\R\times\R^d:\omega=-|\xi|^2\},
 \qquad T_S=\mathcal E_S\mathcal E_S^*.
\]
Frank and Sabin prove mixed-norm Schatten estimates for $q>d+1$ and, in
Remark~5 on PDF page~6 of \cite{FS16}, isolate the remaining endpoint.  In
their notation the proposed estimate is
\begin{equation}\label{eq:source}
 \|\overline W T_S W\|_{\Sp^{d+1,\infty}(L^2(\R^{d+1}))}
 \lesssim
 \|W\|_{L_t^{2(d+1)}L_x^{d+1}}^2 .
\end{equation}
They state that the corresponding strong $\Sp^{d+1}$ estimate is false and
that \eqref{eq:source} was completely open.  The source evidence is reproduced
in \texttt{figures/source\_endpoint\_page6.png}.

\section{The exact dual formulation}

Put $V=|W|^2$ and define the operator on $L^2(\R^d)$
\[
 A_V=\int_{\R} e^{it\Delta}V(t,\cdot)e^{-it\Delta}\,dt .
\]
The usual $AA^*/A^*A$ correspondence identifies the nonzero spectral data in
\eqref{eq:source} with that of $A_V$.  For
\(r=d+1\) and \(r'=(d+1)/d\), Lorentz--Schatten duality gives
\((\Sp^{r',1})^*=\Sp^{r,\infty}\).  Moreover,
\begin{equation}\label{eq:pairing}
 \operatorname{Tr}(\Gamma A_V)
 =\int_{\R}\!\int_{\R^d}
 V(t,x)\,\rho_{e^{-it\Delta}\Gamma e^{it\Delta}}(x)\,dx\,dt .
\end{equation}
Mixed Lebesgue duality therefore makes \eqref{eq:source} equivalent to
\begin{equation}\label{eq:density}
 \left\|\sum_j\lambda_j|e^{it\Delta}f_j|^2\right\|_
 {L_t^{(d+1)/d}L_x^{(d+1)/(d-1)}}
 \lesssim \|\lambda\|_{\ell^{(d+1)/d,1}}
\end{equation}
for orthonormal $(f_j)$.  To pass from the diagonal source formulation to an
arbitrary complex dual test $V$, decompose its real and imaginary parts into
positive and negative parts; each nonnegative part is $|W|^2$.  Since
$d+1>1$, the weak Schatten ideal is normable, so this decomposition loses only
an absolute constant.

For $d=1$, \eqref{eq:source} and \eqref{eq:density} become, respectively,
\[
 \|\overline W T_S W\|_{\Sp^{2,\infty}}
 \lesssim \|W\|_{L_t^4L_x^2}^2,
 \qquad
 \left\|\sum_j\lambda_j|e^{it\Delta}f_j|^2\right\|_{L_t^2L_x^\infty}
 \lesssim \|\lambda\|_{\ell^{2,1}}.
\]

\section{The later counterexample}

Bez, Hong, Lee, Nakamura, and Sawano formulate \eqref{eq:density} as their
Conjecture~1.3.  Their Theorem~5.3 on PDF page~27 states that, when $d=1$,
the induced endpoint velocity-averaging estimate
\begin{equation}\label{eq:kinetic}
 \left\|\int_{\R} f(x-tv,v)\,dv\right\|_{L_t^2L_x^\infty}
 \lesssim \|f\|_{L^{2,1}(\R^2)}
\end{equation}
fails, and hence their conjecture is false \cite{BHLNS17}.  Proposition~5.1
of the same paper proves by semiclassical Weyl quantization that
\eqref{eq:density} would imply \eqref{eq:kinetic}.  The proof of
Theorem~5.3, continued on page~28, tests \eqref{eq:kinetic} on shrinking
neighborhoods of a planar measure-zero Kakeya set.  Thus \eqref{eq:density}
is false for $d=1$, and the duality in \eqref{eq:pairing} proves:

\medskip
\noindent\textbf{Conclusion.}
There is no constant $C$ such that
\[
 \|\overline W T_S W\|_{\Sp^{2,\infty}}
 \le C\|W\|_{L_t^4L_x^2}^2
\]
for every $W$.  This is a full negative answer to the source endpoint in
dimension one.

\section{Scope and current status}

The Kakeya proof uses the endpoint spatial norm $L_x^\infty$, special to
$d=1$.  It does not settle higher dimensions.  Feng, Song, and Wu restate
\eqref{eq:density} as Conjecture~1.6 and, immediately after it on PDF page~6,
write that it is false for $d=1$ but remains open for $d\ge2$ \cite{FSW25}.
Their status page is reproduced in
\texttt{figures/current\_status\_page6.png}.

The literature search checked the exact weak-Schatten and restricted
orthonormal endpoint phrases, the earlier strict-range theorem
arXiv:1404.2817, the 2017 counterexample, and the 2025 critical-summability
paper.  No later all-dimensional proof or counterexample was found.  Eight
focused upgrade routes are recorded separately; none overcame the surviving
mixed-norm/Kakeya endpoint for $d\ge2$.

\begin{thebibliography}{9}
\bibitem{FS16}
R.~L. Frank and J. Sabin,
\emph{The Stein--Tomas inequality in trace ideals},
S\'eminaire Laurent Schwartz---EDP et applications (2015--2016), Exp.~XV;
arXiv:1609.08388.

\bibitem{BHLNS17}
N. Bez, Y. Hong, S. Lee, S. Nakamura, and Y. Sawano,
\emph{On the Strichartz estimates for orthonormal systems of initial data with
regularity}, arXiv:1708.05588, Theorem~5.3.

\bibitem{FSW25}
G. Feng, M. Song, and H. Wu,
\emph{Strichartz estimates involving orthonormal systems at the critical
summability exponent}, arXiv:2507.14974, Conjecture~1.6.
\end{thebibliography}

\end{document}

